\chapter{What Learning Means: Generalisation, Bias--Variance and Regularisation}
\companion{StatQuest: Machine Learning Fundamentals, Bias and Variance}{\ytid{EuBBz3bI-aA}}

Almost every interview question in Rounds 1 and 2 is, underneath, a question about one idea:
\emph{a model must do well on data it has never seen}. This chapter builds that idea from nothing.
Every later chapter (trees, boosting, SVMs, neural networks) is a different answer to the same
problem, so time spent here pays off everywhere.

\section{The problem machine learning solves}

Suppose Naukri wants to show each job seeker an estimate: ``jobs like this usually pay about
Rs.~$X$ lakh''. Nobody can write that rule by hand; it depends on title, city, experience, company,
skills and a hundred other things. But Naukri has \emph{examples}: millions of past jobs with their
salaries. Machine learning is the craft of turning examples into a rule.

We use a few words for this throughout the book.
\begin{itemize}
\item An \textbf{example} (also: row, sample, observation) is one item, such as one job posting.
\item A \textbf{feature} (also: input, variable, column, $x$) is one measured property of it: years of
  experience required, city, number of skills listed.
\item The \textbf{target} (also: label, output, $y$) is what we want to predict: the salary.
\item A \textbf{model} is a rule $f$ that takes features and returns a prediction $\hat y = f(x)$.
  The hat on $\hat y$ always means ``predicted''.
\item \textbf{Parameters} are the numbers inside the model that training adjusts (for example the slope of a
  line). \textbf{Hyperparameters} are the settings \emph{we} choose before training (for example how
  deep a tree may grow).
\item A \textbf{loss} measures how wrong one prediction is. Training means adjusting the parameters to
  make the average loss small on the examples we have.
\end{itemize}

\begin{example}[: the smallest possible model]
Three past jobs: experience $x = 1, 2, 3$ years, salary $y = 4, 6, 8$ lakh. Try the rule ``salary $= 2x + 2$''.
Predictions: $4, 6, 8$. Every error is zero, so the average loss is zero. A new job asking for $x=5$ years gets
the prediction $2\cdot5+2 = 12$ lakh. That last step, predicting for an $x$ we have not seen, is the whole
point. Fitting the three known points was only a means to it.
\end{example}

Two kinds of problem cover most of the syllabus.
\begin{itemize}
\item \textbf{Regression}: the target is a number (salary, property price, days to fill a vacancy).
\item \textbf{Classification}: the target is a category (will the seeker apply: yes/no; which of 500 job
  categories does this title belong to; is this message spam).
\end{itemize}
There is also \textbf{unsupervised} learning, where there is no target at all and we look for structure
(clusters of similar candidates, a compressed representation of resumes). And \textbf{ranking}, where we
order items (which 20 of 30 lakh candidates should a recruiter see first); ranking is central at InfoEdge
because search and recommendation are its core products.

\section{Generalisation: the only thing that matters}

\subsection{Training error versus test error}
The \textbf{training error} is the average loss on the examples used to fit the model. The
\textbf{test error} (or generalisation error) is the average loss on \emph{new} examples from the same
source. We care only about the second; the first is easy to make tiny and tells us little.

\begin{intuition}[: the student who memorised the answer key]
A student who memorises last year's answer key scores 100\% on last year's paper (training error zero) and
fails this year's paper (test error high). A student who understood the chapter scores 85\% on both.
We want the second student. Memorising is called \textbf{overfitting}; understanding is
\textbf{generalising}.
\end{intuition}

\subsection{Why we split data}
Because we cannot see the future, we \emph{simulate} it. We hide part of our data, fit on the rest, and
measure error on the hidden part.
\begin{itemize}
\item \textbf{Training set} (typically 60--80\%): used to fit parameters.
\item \textbf{Validation set} (10--20\%): used to choose hyperparameters and compare models.
\item \textbf{Test set} (10--20\%): touched \emph{once}, at the end, to report an honest number.
\end{itemize}
Why three and not two? Every time you look at a set and make a choice because of it, that set stops being
``unseen''. If you try 50 settings and keep the one with the best validation score, that score is
optimistic: you picked the luckiest of 50. The untouched test set corrects for that.

\begin{example}[: why the best of many is optimistic]
Imagine 50 models that are all \emph{exactly} 80\% accurate in truth. On a validation set of 200 rows,
each one's measured accuracy wobbles by chance, roughly $\pm 2.8\%$ (the standard error
$\sqrt{0.8\cdot0.2/200}\approx 0.028$). The best of 50 such wobbles is usually about 2 standard errors
above the truth, so you would report about 85--86\% for a model that is really 80\%. A fresh test set would
show 80\%.
\end{example}

\subsection{Cross-validation}
When data are scarce, a single validation split is noisy. \textbf{$k$-fold cross-validation} splits the
training data into $k$ equal parts (folds). Train on $k-1$ folds, validate on the remaining one, repeat $k$
times so every fold is the validation fold once, and average the $k$ scores.

\begin{example}[: 5-fold CV on 1{,}000 rows]
Each fold has 200 rows. Round 1: train on rows 201--1000, validate on 1--200. Round 2: train on 1--200 and
401--1000, validate on 201--400, and so on. You get 5 scores, say 0.81, 0.79, 0.83, 0.80, 0.82. Report the
mean 0.81 and the spread (standard deviation about 0.016). Every row was used for validation exactly once and
for training four times.
\end{example}

Variants you should be able to name:
\begin{itemize}
\item \textbf{Stratified $k$-fold}: each fold keeps the class proportions (essential when only 2\% are
  positive; otherwise a fold may contain no positives at all).
\item \textbf{Group $k$-fold}: all rows of one group (one user, one company) stay in the same fold.
  If one recruiter posted 300 jobs, splitting their jobs across train and validation lets the model
  ``recognise the recruiter'' and look better than it is.
\item \textbf{Time-series split}: always train on the past and validate on the future. Random splits on
  time data let the model peek at tomorrow.
\item \textbf{Leave-one-out}: $k = n$. Nearly unbiased but expensive and high-variance.
\end{itemize}

\subsection{Data leakage}
\textbf{Leakage} is when information that would not be available at prediction time sneaks into
training. It produces models that look brilliant offline and fail in production. It is the single most
common reason real projects fail, which is why the InfoEdge coding round rewards spotting it first.

\begin{practical}[: the reported Round 3 leakage]
An MBA placement dataset had a \code{salary} column that was filled only for placed students. If you keep it,
the model learns ``salary present $\Rightarrow$ placed'' and scores near 100\%. At prediction time (before
placement) salary does not exist. The candidate who said ``drop salary, it leaks the target'' before writing
any model code impressed the panel.
\end{practical}

Common leakage patterns:
\begin{itemize}
\item \textbf{Target leakage}: a feature that is a consequence of the target (salary after placement;
  ``number of interview calls'' when predicting ``got the job'').
\item \textbf{Preprocessing leakage}: computing the mean for imputation, the scaler, TF-IDF vocabulary or
  PCA on the \emph{whole} dataset before splitting. The validation rows then influenced training.
\item \textbf{Temporal leakage}: using features computed with future data (``total applications this job
  will receive'').
\item \textbf{Duplicate leakage}: the same resume uploaded twice landing in both train and test.
\end{itemize}

\section{Underfitting, overfitting and model capacity}

\textbf{Capacity} (or complexity) is how flexible a model is: how many different shapes it can bend into.
A straight line has low capacity; a degree-15 polynomial or a deep tree has high capacity.

\begin{center}
\begin{tikzpicture}
\begin{axis}[width=0.62\linewidth,height=5.2cm,xlabel={model complexity},ylabel={error},
  xmin=0,xmax=10,ymin=0,ymax=10,xtick=\empty,ytick=\empty,legend pos=north east,legend style={font=\small}]
\addplot[domain=0.5:10,samples=60,thick,accent]{8*exp(-0.45*x)+0.3};
\addlegendentry{training error}
\addplot[domain=0.5:10,samples=60,thick,prframe]{8*exp(-0.45*x)+0.06*x^2+1.0};
\addlegendentry{test error}
\draw[dashed] (axis cs:4.3,0) -- (axis cs:4.3,10);
\node[font=\small] at (axis cs:1.8,9.3) {underfit};
\node[font=\small] at (axis cs:8.3,9.3) {overfit};
\node[font=\small] at (axis cs:4.3,0.6) {sweet spot};
\end{axis}
\end{tikzpicture}
\end{center}

\begin{itemize}
\item \textbf{Underfitting}: training error is high and test error is about as high. The model is too
  simple to capture the pattern (a straight line through data that curve).
\item \textbf{Overfitting}: training error is low but test error is much higher. The model learned the
  noise in the particular sample.
\item Training error falls steadily as complexity rises. Test error first falls, then rises: a U shape.
  The bottom of the U is what we tune for.
\end{itemize}

\begin{example}[: reading the gap]
\begin{tabular}{lccc}
Model & Train acc. & Val. acc. & Diagnosis\\\hline
A & 62\% & 61\% & underfit (both low, small gap)\\
B & 99\% & 72\% & overfit (big gap)\\
C & 88\% & 86\% & good (high and small gap)\\
\end{tabular}
\end{example}

\section{The bias--variance decomposition, built slowly}

\subsection{An experiment in your head}
Imagine you could collect \emph{many} different training sets of the same size from the same world, and
train the same kind of model on each. For one fixed test input $x$ (say a job with 5 years' experience in
Pune), you would get many different predictions, one per training set. Look at that cloud of predictions.
\begin{itemize}
\item \textbf{Bias} is how far the \emph{centre} of the cloud is from the truth. A model that is
  systematically wrong (a straight line for a curved truth) has high bias, however much data it gets.
\item \textbf{Variance} is how \emph{spread out} the cloud is. A model that changes a lot when a few
  training rows change (a deep tree) has high variance.
\item \textbf{Noise} is the randomness in $y$ itself that no model can predict: two identical jobs can
  pay differently because of negotiation.
\end{itemize}

\begin{intuition}[: the dartboard]
Four archers. Low bias, low variance: tight cluster on the bullseye. High bias, low variance: tight cluster,
but off-centre (consistently wrong). Low bias, high variance: centred on average but scattered. High bias,
high variance: scattered and off-centre. Training many models on many datasets is throwing many darts.
\end{intuition}

\subsection{The formula, and what each piece means}
For squared error, the expected test error at a point $x$ splits exactly into three parts:
\[
\E\big[(y - \hat f(x))^2\big] \;=\; \underbrace{\big(\E[\hat f(x)] - f(x)\big)^2}_{\text{bias}^2}
\;+\; \underbrace{\E\big[(\hat f(x) - \E[\hat f(x)])^2\big]}_{\text{variance}}
\;+\; \underbrace{\sigma^2}_{\text{noise}}.
\]
Here $f(x)$ is the true average target at $x$, $\hat f(x)$ is our model's prediction (random, because it
depends on the random training set), $\E[\cdot]$ means ``average over many training sets'', and $\sigma^2$
is the noise variance.

\begin{example}[: computing the three parts]
At $x$ = ``5 years, Pune'' the true average salary is $f(x) = 10$ lakh and the noise variance is
$\sigma^2 = 1$. Over many training sets our model predicts $9, 11, 13, 11$ (four sets, for simplicity).
Mean prediction $= 11$, so bias $= 11 - 10 = 1$ and bias$^2 = 1$. Variance $= \frac{(9-11)^2 + 0 + (13-11)^2 + 0}{4}
= \frac{8}{4} = 2$. Expected error $= 1 + 2 + 1 = 4$, i.e. a typical miss of $\sqrt4 = 2$ lakh.
\end{example}

\subsection{Why there is a trade-off}
Making a model more flexible lets its average move closer to the truth (bias falls) but also lets it chase
the quirks of each sample (variance rises). Making it more rigid does the opposite. Total error is minimised
somewhere in between. More \emph{data} is the one lever that lowers variance without raising bias, which is
why ``get more data'' is the right answer to an overfitting model and the wrong answer to an underfitting one.

\begin{center}\small
\begin{tabularx}{\linewidth}{l X X}
\toprule
& High bias (underfit) & High variance (overfit)\\
\midrule
Symptom & train and val error both high, close together & train error low, val error much higher\\
Typical models & linear model on curved data; very shallow tree; $k$-NN with huge $k$ & deep tree; $k$-NN with
$k=1$; huge network on little data\\
Fixes & more features, more flexible model, less regularisation, train longer &
more data, regularisation, simpler model, bagging, dropout, early stopping, data augmentation\\
\bottomrule
\end{tabularx}
\end{center}

\subsection{Learning curves: the diagnostic picture}
Plot training and validation error against the \emph{number of training examples}.
\begin{itemize}
\item High bias: both curves flatten quickly at a high error and nearly meet. More data will not help.
\item High variance: a large gap that narrows slowly as data grow. More data will help.
\end{itemize}
A second curve InfoEdge reportedly asks about is error against \emph{training iterations} for a neural
network: training error keeps falling; validation error falls, bottoms out, then rises. Stopping at the
bottom is \textbf{early stopping}.

\subsection{A modern footnote: double descent}
Very large models (modern deep nets) can show a second descent: test error rises near the point where the
model can just fit the training data exactly, then \emph{falls again} as the model grows far beyond it. The
classical U-curve is still the right first answer in an interview; mention double descent as a nuance if
asked about deep learning.

\section{Regularisation: paying for complexity}

\subsection{The idea}
\textbf{Regularisation} is anything that discourages a model from becoming more complex than the data
justify. The most common form adds a penalty to the loss:
\[
\text{minimise}\quad \underbrace{\text{average loss on training data}}_{\text{fit the data}}
\;+\; \lambda \cdot \underbrace{\text{complexity}(\text{parameters})}_{\text{stay simple}}.
\]
$\lambda \ge 0$ (lambda) is a hyperparameter: $\lambda = 0$ means no penalty; large $\lambda$ means simplicity
matters a lot.

\begin{intuition}[: the budget]
Think of every weight as costing money. A feature must ``earn'' its weight by reducing the loss more than it
costs. Weak, noisy features cannot afford a large weight, so their influence shrinks. That is exactly how
regularisation cuts variance.
\end{intuition}

\subsection{L2 (Ridge) and L1 (Lasso)}
For a model with weights $w_1,\dots,w_d$:
\begin{itemize}
\item \textbf{L2 / Ridge / weight decay:} penalty $\lambda \sum_j w_j^2$. Shrinks all weights smoothly toward
  zero but rarely makes any exactly zero.
\item \textbf{L1 / Lasso:} penalty $\lambda \sum_j |w_j|$. Pushes many weights to \emph{exactly} zero, so it
  also selects features.
\item \textbf{Elastic net:} both. Useful when features come in correlated groups (Lasso alone picks one of
  a group arbitrarily).
\end{itemize}

\begin{example}[: one weight, by hand]
Suppose without a penalty the best weight would be $w = 3$, and the loss near it is $\frac12(w-3)^2$.
\begin{itemize}
\item Ridge with penalty $\frac{\lambda}{2}w^2$, $\lambda = 1$: minimise $\frac12(w-3)^2 + \frac12 w^2$. Derivative
  $(w-3) + w = 0 \Rightarrow w = 1.5$. Shrunk by half, but not zero.
\item Lasso with penalty $\lambda|w|$, $\lambda = 1$: for $w>0$ the derivative is $(w-3) + 1 = 0 \Rightarrow w = 2$.
  If instead the unpenalised optimum were $w=0.8$ (a weak feature), the same calculation gives $w = -0.2$,
  which contradicts $w>0$; the minimum is at exactly $w=0$. Lasso subtracts a fixed amount $\lambda$ and stops at
  zero (``soft thresholding''): $w = \text{sign}(w_0)\max(|w_0| - \lambda, 0)$.
\end{itemize}
\end{example}

\subsubsection{Why L1 gives exact zeros: the picture}
Minimising loss $+\lambda\cdot$penalty is equivalent to minimising the loss subject to a budget
``penalty $\le t$''. In two dimensions the L2 budget $w_1^2 + w_2^2 \le t$ is a circle; the L1 budget
$|w_1|+|w_2| \le t$ is a diamond with corners on the axes. The loss contours are ellipses centred at the
unpenalised solution; we expand them until they first touch the budget region. A circle is smooth, so the
touching point is almost never on an axis. A diamond has sharp corners sticking out, and an expanding
ellipse very often hits a corner first. A corner lies on an axis, so one weight is exactly zero.

\begin{center}
\begin{tikzpicture}[scale=0.85]
\begin{scope}
\draw[->] (-2,0)--(2.4,0) node[right]{$w_1$}; \draw[->] (0,-2)--(0,2.4) node[above]{$w_2$};
\draw[thick,accent] (0,0) circle (1);
\draw[prframe] (1.6,1.7) ellipse (1.25 and 0.6);
\draw[prframe] (1.6,1.7) ellipse (0.8 and 0.35);
\fill (0.72,0.69) circle (1.5pt);
\node[font=\small] at (0,-2.4) {L2: touches off-axis};
\end{scope}
\begin{scope}[xshift=6cm]
\draw[->] (-2,0)--(2.4,0) node[right]{$w_1$}; \draw[->] (0,-2)--(0,2.4) node[above]{$w_2$};
\draw[thick,accent] (1,0)--(0,1)--(-1,0)--(0,-1)--cycle;
\draw[prframe,rotate around={-25:(1.2,1.9)}] (1.2,1.9) ellipse (1.25 and 0.55);
\draw[prframe,rotate around={-25:(1.2,1.9)}] (1.2,1.9) ellipse (0.7 and 0.3);
\fill (0,1) circle (1.5pt);
\node[font=\small] at (0,-2.4) {L1: touches a corner, $w_1=0$};
\end{scope}
\end{tikzpicture}
\end{center}

\subsubsection{The Bayesian view (one sentence, often asked)}
L2 regularisation is the same as assuming a Gaussian prior on the weights (most weights small, centred at
zero); L1 is the same as a Laplace prior (sharply peaked at zero). Chapter 12 derives this.

\subsection{Other regularisers you must be able to name}
\begin{itemize}
\item \textbf{Early stopping}: stop training when validation error starts rising.
\item \textbf{Dropout} (neural nets): randomly switch off units during training.
\item \textbf{Data augmentation}: create plausible new training examples (flipped images, paraphrased text).
\item \textbf{Tree limits}: maximum depth, minimum samples per leaf, pruning.
\item \textbf{Ensembling} (bagging) reduces variance by averaging.
\item \textbf{Label smoothing}, \textbf{weight sharing} (CNNs), \textbf{batch norm} (a mild side effect).
\end{itemize}

\begin{trap}
Always \textbf{standardise features} before Ridge or Lasso. The penalty treats every weight equally, but a
weight's size depends on the feature's units: ``salary in rupees'' needs a tiny weight, ``salary in crores''
a huge one. Without scaling, the penalty unfairly punishes features measured in large-number units less and
small-unit features more.
\end{trap}

\section{No free lunch, and why baselines matter}
The \textbf{no free lunch theorem} says that averaged over \emph{all} possible problems, no learning algorithm
beats any other. Every algorithm wins by making assumptions (smoothness, linearity, locality) that suit the
problem. Practically: always start with a simple \textbf{baseline} (predict the mean; logistic regression;
``most popular jobs'') and demand that a complex model beat it by a margin worth its cost.

\stuck{StatQuest: Machine Learning Fundamentals, Cross Validation}{\ytid{fSytzGwwBVw}}
\section{Interview questions}

\begin{qa}{\pyq{OA and INT: bias--variance graphs} Draw the error versus model complexity graph and
explain it. Where is the best model?}
\oneline{Training error falls monotonically with complexity; test error is U-shaped; the best model sits at
the bottom of the test-error U, where the sum of bias$^2$ and variance is smallest.}
\why
On the x-axis is model complexity (depth of a tree, degree of a polynomial, number of parameters, or
$1/\lambda$). On the y-axis, error.
\begin{enumerate}
\item The \textbf{training error} curve starts high (a very simple model cannot fit even the data it sees)
  and keeps falling. A sufficiently complex model can memorise every training point, so it can reach zero.
\item The \textbf{test error} curve first falls together with training error, because extra flexibility
  is capturing real structure (bias falling). Beyond some point, extra flexibility only captures noise
  specific to the sample (variance rising), so test error turns upward.
\item The left region is \textbf{underfitting} (high bias), the right region \textbf{overfitting} (high
  variance). The vertical gap between the curves is a measure of overfitting.
\end{enumerate}
The best model is at the minimum of the \emph{validation} curve (we estimate the test curve with a validation
set or cross-validation).
\worked A polynomial fitted to 20 noisy points from a sine curve: degree 1 has train MSE 0.30, val 0.32
(underfit); degree 4 has 0.05 and 0.07 (good); degree 15 has 0.001 and 2.4 (overfit). Choose degree 4.
\pushes
\emph{``What if I plot against number of training examples instead?''} That is a learning curve: training
error rises (harder to fit more points), validation error falls; they converge. A persistent big gap
signals variance; both converging at a high value signals bias.
\emph{``And against training iterations of a neural net?''} Training error keeps falling; validation error
falls then rises: stop at its minimum (early stopping).
\pitfall Saying ``test error decreases forever with complexity'' or putting the optimum at zero training
error. Also do not confuse the x-axis: more \emph{data} is not more \emph{complexity}.
\end{qa}

\begin{qa}{\pyq{OA, IIT Roorkee ML section} A model has 99\% training accuracy and 72\% validation accuracy.
What is happening, and what would you try, in order?}
\oneline{It is overfitting (high variance); I would first check for leakage or a split problem, then add
regularisation or reduce complexity, and get more or augmented data if possible.}
\why
A 27-point gap means the model's behaviour depends heavily on which rows it saw. Before changing the model I
verify the measurement:
\begin{enumerate}
\item \textbf{Is the split sound?} Are there duplicates across train and validation? Is the validation set
  from a different period or distribution (a time shift)? Is the validation set tiny and noisy?
\item \textbf{Is there leakage in training only?} For example a feature that exists in the training extract
  but is computed differently for validation.
\end{enumerate}
If the measurement is fine, reduce variance:
\begin{enumerate}
\item Simplify: shallower trees, fewer features, smaller network.
\item Regularise: L1/L2, dropout, min-samples-per-leaf, early stopping.
\item Average: bagging or a random forest instead of one tree.
\item More data or augmentation: the only fix that reduces variance without adding bias.
\item Feature selection: remove noisy features (high-cardinality IDs are classic culprits).
\end{enumerate}
\worked A decision tree on 5{,}000 Shiksha leads to predict ``will apply to a college'': unlimited depth gives
99\%/72\%. Limiting \code{max\_depth=6} and \code{min\_samples\_leaf=50} gives 84\%/81\%. A random forest
gives 90\%/85\%.
\pushes \emph{``Why not just collect more data first?''} It is often slow and expensive; regularisation is
free. But if a learning curve shows validation error still dropping steeply with more data, data collection
is worth it.
\pitfall Jumping straight to ``use a bigger model'' (that worsens variance), or to SMOTE, which addresses class
imbalance, not overfitting.
\end{qa}

\begin{qa}{\pyq{INT: weight decay} What is weight decay? Is it the same as L2 regularisation?}
\oneline{Weight decay shrinks every weight by a small fraction at each update; for plain gradient descent it is
exactly equivalent to adding an L2 penalty, but for adaptive optimisers like Adam the two differ, which is why
AdamW exists.}
\why
With L2, the loss becomes $L(w) + \frac{\lambda}{2}\norm{w}^2$, whose gradient is $\nabla L + \lambda w$. A
gradient-descent step is
\[ w \leftarrow w - \eta(\nabla L + \lambda w) = (1-\eta\lambda)\,w - \eta \nabla L. \]
The factor $(1-\eta\lambda)$ multiplies (``decays'') the weight every step before the usual update. That is
weight decay. So for SGD, ``weight decay with coefficient $\eta\lambda$'' and ``L2 with coefficient $\lambda$''
are the same thing.

In Adam, the gradient is divided by a running estimate of its size, per weight. If the L2 term is put
inside the gradient, it also gets divided, so weights with large gradients are decayed \emph{less}. That is not
what we want. \textbf{AdamW} applies the decay directly to the weights, outside the adaptive scaling
(``decoupled weight decay''), and generalises better; it is the default for training Transformers.
\worked $w = 1$, $\eta = 0.1$, $\lambda = 0.01$, and the data gradient happens to be 0. One SGD step: $w \leftarrow
(1 - 0.001)\cdot 1 = 0.999$. With no data pushing it, the weight decays geometrically toward 0.
\pushes \emph{``Why does shrinking weights reduce overfitting?''} Small weights make the function smoother: a
small change in input causes a small change in output, so the model cannot carve sharp, data-specific
wiggles. In a sigmoid/tanh network, small weights also keep units in their near-linear range.
\pitfall Saying weight decay ``reduces the learning rate''. It shrinks the weights, not the step size.
\end{qa}

\begin{qa}{\pyq{INT and OA DL: overfitting} List the ways to reduce overfitting in a neural network and
explain why each works.}
\oneline{More data or augmentation, smaller network, L2/weight decay, dropout, early stopping, batch norm,
label smoothing, and ensembling; each either limits the effective capacity or makes the training signal less
specific to the sample.}
\why
\begin{itemize}
\item \textbf{More data / augmentation}: noise patterns stop being repeatable, so memorising them stops paying.
  Augmentation (crops, flips, synonym replacement) teaches invariances we know are true.
\item \textbf{Smaller network}: fewer parameters, fewer ways to memorise.
\item \textbf{Weight decay}: smooth functions (previous question).
\item \textbf{Dropout}: randomly zeroes units each step with probability $p$, so no unit can rely on a
  specific partner; it behaves like averaging many thinned networks (an implicit ensemble).
\item \textbf{Early stopping}: weights start small and grow during training; stopping early limits how
  far they move, acting like an L2 penalty.
\item \textbf{Batch normalisation}: the batch statistics add a little noise per example, a mild regulariser.
\item \textbf{Label smoothing}: replaces hard 0/1 targets with 0.05/0.95-style targets, so the network is not
  rewarded for infinitely confident logits.
\item \textbf{Ensembling / checkpoint averaging}: averaging reduces variance.
\end{itemize}
\worked A CNN classifying 2{,}000 property photos into 6 room types: 99\% train, 70\% val. Adding flips/crops
and dropout 0.3 gives 93\%/82\%; switching to a pretrained ResNet fine-tuned with weight decay gives 95\%/91\%.
Transfer learning is often the biggest single fix when data are small.
\pitfall Listing techniques without the ``why''. The panel reported asking exactly this to test intuition.
\end{qa}

\begin{qa}{Explain the bias--variance trade-off to a non-technical product manager in under a minute.}
\oneline{A model can be wrong because it is too rigid to see the pattern (bias) or because it is so eager that
it mistakes coincidences for patterns (variance); we tune its flexibility to balance the two.}
\why
``Imagine predicting the price of flats in Noida. If I only use `price = average price', I am always off for
big or small flats: that is \emph{bias}, a systematic blind spot. If I memorise every flat sold last month,
including the one that sold cheap because the owner was in a hurry, my predictions jump around depending on
which flats happened to sell: that is \emph{variance}. The sweet spot is a model that uses size, location and
floor but ignores one-off stories. We find it by testing on flats the model has not seen.''
\pushes \emph{``So which one is worse?''} Neither in general; the business question decides. For a salary
estimate shown to users, high variance (estimates that jump week to week) erodes trust, so we would accept
slightly more bias for stability.
\end{qa}

\begin{qa}{Compute the expected squared error at a point where the true value is 10, predictions over
four training sets are 8, 12, 14, 10, and the noise variance is 2.}
\oneline{Bias$^2$ = 1, variance = 5, noise = 2, so the expected error is 8.}
\why Mean prediction $= (8+12+14+10)/4 = 11$. Bias $= 11 - 10 = 1$, bias$^2 = 1$.
Variance $= \frac{(8-11)^2 + (12-11)^2 + (14-11)^2 + (10-11)^2}{4} = \frac{9+1+9+1}{4} = 5$.
Total $= 1 + 5 + 2 = 8$.
\pushes \emph{``Which part can no model remove?''} The noise, 2. \emph{``How would bagging change this?''} It
averages models, cutting the variance term (towards $\rho\sigma^2$ if models are correlated) while leaving bias
roughly unchanged.
\end{qa}

\begin{qa}{Why does L1 regularisation produce sparse weights but L2 does not? Answer both geometrically and
with derivatives.}
\oneline{The L1 penalty has a constant-size pull toward zero that does not fade as the weight shrinks, and its
constraint region has corners on the axes; the L2 pull is proportional to the weight, so it fades to nothing
near zero.}
\why
\textbf{Derivatives.} The L2 penalty $\lambda w^2$ has derivative $2\lambda w$: when $w$ is already small, the
push toward zero is tiny, so $w$ settles at a small non-zero value wherever the data's pull balances it. The L1
penalty $\lambda|w|$ has derivative $\lambda\cdot\text{sign}(w)$: a constant push of size $\lambda$ however small
$w$ is. If the data's pull on a weight is weaker than $\lambda$, nothing can hold the weight away from zero, so it
lands exactly at zero.

\textbf{Geometry.} As drawn in Section 1.5, the L1 constraint region is a diamond whose corners lie on the
axes; expanding loss ellipses usually touch a corner first, and corners have some coordinates equal to zero.
The L2 region is a circle with no corners.
\worked With an unpenalised optimum $w_0 = 0.8$ and $\lambda = 1$: Lasso gives $\max(0.8 - 1, 0) = 0$ (feature
removed). Ridge with the same loss gives $0.8/(1+1) = 0.4$ (shrunk, kept).
\pushes \emph{``When would you prefer Ridge anyway?''} When many features each carry a little signal (text
features, embeddings) or features are highly correlated: Lasso arbitrarily keeps one of a correlated group and
can be unstable; Ridge spreads weight across the group.
\end{qa}

\begin{qa}{What is data leakage? Give three concrete examples from a job-portal setting and how you would
detect each.}
\oneline{Leakage is using, during training, information that will not be available when the model makes real
predictions; it inflates offline scores and collapses in production.}
\why
\begin{enumerate}
\item \textbf{Target leakage}: predicting ``will this candidate be hired'' using ``number of offer letters
  generated''. Detection: a single feature with almost perfect importance or AUC; ask ``when is this value
  created, relative to the event?''
\item \textbf{Preprocessing leakage}: fitting a scaler or TF-IDF on all resumes before splitting.
  Detection: code review; enforce a \code{Pipeline} so every transform is fit inside each training fold.
\item \textbf{Temporal leakage}: random train/test split on applications from 2023--2025 when the model
  will predict 2026 behaviour; also features like ``total applies this job will get''. Detection: compare a
  random split score with a time-based split score; a large drop signals leakage.
\item \textbf{Group leakage}: the same candidate's multiple applications in both train and test.
  Detection: group $k$-fold by candidate ID.
\end{enumerate}
\pushes \emph{``Your offline AUC is 0.99 on a hard problem. What do you think first?''} Leakage, before
celebrating. Then I look at the top features and the split.
\end{qa}

\begin{qa}{Why do we need a separate test set if we already use cross-validation?}
\oneline{Because choosing the best configuration using CV scores makes those scores optimistic; an untouched
test set gives an unbiased final estimate.}
\why CV is used to \emph{select} among hyperparameters and models. Selection on noisy scores favours lucky
configurations (the maximum of noisy numbers is biased upward). The test set has not influenced any decision,
so its score is an honest estimate of performance on new data. When data are very small, \textbf{nested
cross-validation} replaces the single test set: an outer CV loop estimates performance, an inner CV loop tunes
hyperparameters inside each outer training fold.
\end{qa}

\begin{qa}{When is a random train--test split wrong? Give two cases and the right split.}
\oneline{When rows are not independent: time-ordered data (use a time-based split) and grouped data such as
many rows per user (use a group split).}
\why
A random split assumes rows are exchangeable, i.e. the test rows look like unseen future rows. For time series
(daily applications to a job), a random split puts tomorrow in training and yesterday in testing: the model
interpolates rather than forecasts. Use a forward-chaining split. For grouped data (many sessions per seeker),
a random split lets the model memorise each seeker's habits and get them ``right'' in test. Use
\code{GroupKFold}. A third case is \textbf{distribution shift}: if the model will be deployed in a new city,
validate on a held-out city.
\end{qa}

\begin{qa}{\deep Can adding a completely useless random feature ever \emph{improve} validation accuracy? Can it
hurt?}
\oneline{It can appear to help by chance on a small validation set, but in expectation it hurts slightly,
because it adds variance without adding signal; how much it hurts depends on the model's regularisation.}
\why
A random feature has no relationship with the target in the population. A flexible model will still find
spurious patterns in it on the training sample (pure noise always correlates a little with anything in a
finite sample). Those patterns do not repeat on new data, so predictions get noisier: variance goes up. With
strong regularisation (Lasso, a random forest that rarely picks it) the damage is small; with an unregularised
model on few rows and many such features, it can be large. On a single small validation set, random chance can
make the noisy model score higher, which is exactly why we use CV and look at the spread.
\worked Linear regression with $n = 30$ rows: adding 20 random columns raises training $R^2$ from 0.60 to 0.85
(it can only go up) while CV $R^2$ drops from 0.55 to 0.20.
\pushes \emph{``How is this used in practice?''} A classic feature-selection trick: add a few random ``shadow''
features (Boruta) and drop any real feature whose importance is not higher than the best random one.
\end{qa}

\begin{qa}{\deep A model has \emph{lower} validation loss than training loss. Is something wrong?}
\oneline{Not necessarily; regularisation that is active only during training (dropout, augmentation), the timing
of when training loss is averaged, or an easier validation set can all cause it, but it is worth checking for a
leaky or unrepresentative validation set.}
\why
\begin{enumerate}
\item \textbf{Dropout and augmentation} make training harder: training loss is computed on a crippled
  network and distorted inputs; validation uses the full network on clean inputs.
\item \textbf{Averaging timing}: training loss is often the average over the epoch while the weights were
  still improving; validation loss is computed at the end of the epoch with the better weights.
\item \textbf{Easier validation data}: by chance, or because of a bad split (validation drawn from a cleaner
  source), or because of duplicates of training rows in validation (leakage).
\end{enumerate}
I would check the third explanation first, since the first two are benign.
\end{qa}

\begin{qa}{What is the difference between parameters and hyperparameters? Give three examples of each for
three different models.}
\oneline{Parameters are learned from data by the training algorithm; hyperparameters are chosen by us before
training and control how learning happens.}
\why
\begin{center}\small
\begin{tabular}{lll}
\toprule
Model & Parameters (learned) & Hyperparameters (chosen)\\
\midrule
Linear/logistic regression & weights, bias & $\lambda$ (or $C$), penalty type\\
Decision tree & split features, thresholds, leaf values & max depth, min samples per leaf, criterion\\
Neural network & weights and biases of every layer & learning rate, layers, width, dropout, batch size\\
$k$-means & centroid positions & $k$, initialisation\\
\bottomrule
\end{tabular}
\end{center}
Hyperparameters are tuned on validation data (grid, random or Bayesian search), never on the test set.
\end{qa}

\begin{qa}{Your model has 62\% training and 61\% validation accuracy on a problem where humans get 95\%. What
do you do?}
\oneline{This is high bias (underfitting): increase capacity or improve features, and reduce regularisation;
more data alone will not help.}
\why The gap is tiny, so the model is consistent but systematically limited. Options in order:
(1) richer features (interactions, text embeddings, domain features); (2) a more flexible model (gradient
boosting instead of a linear model; a deeper network); (3) less regularisation (larger $C$, smaller
$\lambda$, less dropout); (4) train longer if it is a neural network that has not converged;
(5) check label quality: if labels are noisy, ``human 95\%'' may not be achievable with these labels.
\pitfall Recommending more data. A learning curve that is flat at 61\% tells you more rows will not move it.
\end{qa}

\begin{qa}{\deep Is a model with 0\% training error always overfitting?}
\oneline{No. Zero training error means the model can fit the sample; it is overfitting only if validation error is
much worse than it could be. Many good models (boosted trees, large neural nets, 1-NN on a clean problem)
reach zero training error and still generalise well.}
\why Overfitting is defined by the \emph{gap and by comparison with alternatives}, not by training error alone.
If the classes are perfectly separable and labels are clean (say, digit images), fitting every training point
can be correct. Modern deep networks routinely interpolate the training set yet generalise, which is part of the
double-descent story. Conversely, if labels are noisy (10\% mislabelled), zero training error means the model
learned the mislabelled points too, which usually hurts.
\end{qa}

\begin{qa}{What is $k$-fold cross-validation, how do you choose $k$, and what does it cost?}
\oneline{Split data into $k$ folds, train $k$ times each time holding one fold out, and average the scores;
$k = 5$ or $10$ is the usual compromise between bias of the estimate, its variance, and compute.}
\why Small $k$ (like 2) trains on only half the data each time, so each model is weaker than the final one
(pessimistic, biased estimate). Large $k$ (like $n$, leave-one-out) trains on almost all data (low bias) but the
$k$ models are nearly identical, so their errors are correlated and the estimate is noisy; it also costs $k$ full
trainings. With 1 million rows, a single holdout is already precise, so CV is often unnecessary; with 2{,}000
rows, 5- or 10-fold stratified CV is standard.
\pushes \emph{``After CV, which model do you deploy?''} Usually retrain on all training data with the chosen
hyperparameters (or average the $k$ fold models, common in competitions).
\end{qa}

\begin{qa}{\deep If regularisation always increases training error, why does it ever help?}
\oneline{Because it trades a small increase in bias for a larger decrease in variance, and test error is the sum
of both.}
\why Training error measures fit to this sample. Regularisation deliberately stops the model from fitting the
sample perfectly, so training error rises. But the parts of the sample it refuses to fit are disproportionately
noise. On new data that noise is different, so not having fitted it is an advantage. In bias--variance terms:
variance falls faster than bias$^2$ rises, until $\lambda$ becomes so large that the model ignores real signal
too, after which test error rises again. That is why $\lambda$ is tuned on validation data.
\worked Ridge on 50 rows, 40 features: $\lambda = 0$: train MSE 0.1, val 5.0. $\lambda = 1$: train 0.8, val 1.2.
$\lambda = 100$: train 3.0, val 3.1 (now underfitting).
\end{qa}

\begin{qa}{What is the no-free-lunch theorem and what does it mean for a working data scientist?}
\oneline{No algorithm is best on every possible problem; an algorithm wins only when its built-in assumptions
match the data, so we must try sensible candidates and validate empirically.}
\why A linear model assumes additive effects; a CNN assumes local, translation-equivariant patterns; $k$-NN
assumes nearby points share labels; gradient-boosted trees assume axis-aligned, piecewise-constant structure.
Each is excellent when its assumption holds and poor otherwise. Practically: (1) start with a baseline,
(2) pick candidates whose assumptions fit the data type (trees for tabular, CNNs for images, Transformers for
text), (3) compare on a sound validation scheme.
\end{qa}

\begin{qa}{\deep Two models have the same validation accuracy, 85\%. One is logistic regression, the other a
2{,}000-tree gradient-boosted model. Which do you ship?}
\oneline{Usually the simpler one: equal accuracy plus lower latency, easier debugging, better calibration and
interpretability, and typically more stable behaviour under data shift.}
\why Accuracy is one of several criteria. Ask: (1) Is the 85\% estimate equally certain for both? Look at CV
spread. (2) Latency and cost: a model scoring 30 lakh Resdex candidates per search must be fast. (3)
Interpretability: recruiters and compliance may need reasons. (4) Robustness: the simpler model has fewer
ways to break when inputs drift. (5) Other metrics: compare calibration, recall at the operating point, and
performance on important segments (freshers, tier-2 cities). If the complex model wins clearly on a metric
that matters (say recall on rare high-value leads), it can be worth it.
\end{qa}

\begin{qa}{What is stratified sampling and when must you use it?}
\oneline{Splitting so that each subset keeps the same class proportions as the whole; it is essential for
imbalanced classification and small datasets.}
\why With 1\% positives and 1{,}000 rows, a random 20\% validation split contains 2 positives on average and
quite often 0 or 1, making recall meaningless and unstable. Stratification guarantees exactly 2 per fold.
In scikit-learn: \code{train\_test\_split(..., stratify=y)} and \code{StratifiedKFold}.
\end{qa}
